\chapter{Introduction}

\DepPy{} is easiest to misunderstand when its different subsystems are read as separate
one-off experiments: a plain type hierarchy here, a hybrid dependent type story there, a
cohomological bug-finding paper somewhere else, and an agent-driven generation runtime on
top. The working thesis of this manuscript is that those pieces are better read as one
architecture. The architecture begins from a simple observation: software understanding is
local-to-global. We learn facts at boundaries, guards, call sites, proof obligations,
library interfaces, and traces. We then try to compose those local facts into a global
judgment about safety, meaning, equivalence, or implementation quality. When local facts
agree, we obtain a global section. When they fail to agree, we do not merely have a vague
error---we have a gluing obstruction.

That is why cohomological language keeps reappearing across the repository. In the older
sheaf-oriented work, a program is analyzed through a site category of observation points,
a semantic presheaf of local facts, and the computation of \v{C}ech cohomology to locate
where those facts fail to glue. In the hybrid work, the same local-to-global idea is
lifted to a five-layer product site where intent, structural facts, semantic judgments,
proof artifacts, and code must agree. In the equivalence machinery, one asks whether
local isomorphisms glue to a global one. In anti-hallucination, one asks whether a
candidate artifact inhabits the expected type on both its structural and semantic faces.
In prompt-driven generation, one grows a cover of typed module obligations and tries to
glue them into a coherent project.

The phrase ``cohomological typing for absolutely everything'' is intentionally ambitious,
but it is not meant as parody. It means that the same mathematical stance can organize
far more of the repository than ordinary type-checking alone. Typing becomes a universal
language of structured local evidence: value-level facts, control-flow facts, library
facts, semantic judgments, proof obligations, and design intent all enter the same story.
The right reaction is not to believe that every workflow literally computes a heavy
cohomology group in the same way. The right reaction is to notice that the repository has
already started treating the local-to-global pattern as its common semantic denominator.
This monograph makes that denominator explicit.

It is worth saying more about what counts as a single project here. A manuscript can become
long without becoming unified. It can merely accumulate examples, slogans, and modules. A
foundational project earns unity only when its later constructions seem to have been latent
in the earlier ones all along. The burden of this monograph is therefore not simply to
explain many parts of the repository. It is to persuade the reader that once one grants the
local-to-global viewpoint, the rest of the repository becomes a chain of increasingly rich
consequences: first exact obligations, then layered evidence, then cohomological diagnosis,
then trust-sensitive repair, then exploratory elaboration, and finally agent-driven
assembly. The paper should feel like a single river that keeps widening, not like a museum
of related specimens.

The two most important distinctions in what follows are these. First, the plain type
system in \repo{src/deppy/types/} is not superseded by the hybrid system; it is the
object language on which the hybrid system depends. It provides the type constructors,
refinement syntax, binder operations, and subtype structure that later layers continue to
use. Second, the hybrid system in \repo{src/deppy/hybrid/} is not merely ``the theorem
prover side'' or ``the LLM side.'' It is a broader extension that adds verification
layers, trust, provenance, semantic predicates, Lean translation, anti-hallucination,
zero-change adoption, and prompt-driven planning. The hybrid story is therefore both more
mathematical and more practical than a simple proof layer.

\paragraph{What this document is trying to do.}
The document has four goals. First, it gives a coherent account of the foundational plain
type system and the layered hybrid extension. Second, it explains how site categories,
presheaves, trust lattices, and dependent formers organize the repository's major ideas.
Third, it shows how apparently distant components---theory packs, renderers, equivalence,
ideation, generation agents, and trading examples---fit into the same architecture.
Fourth, it tries to be honest about tensions in the repository: duplicated trust vocabularies,
heuristic bridges from intent to formalization, and the difference between semantic
acceptance and formal theorem proving.

\paragraph{Why the comparison to a foundational project matters.}
The aspiration behind the present document is not to mimic the vocabulary of better-known
foundational programs, but to borrow one of their virtues: they make the reader feel that a
few basic ideas change the subject all the way down. In that spirit, the early chapters here
will insist on sites, sections, and exact obligations; the middle chapters will insist on
layered evidence, trust, and cohomological obstruction; the later chapters will insist that
generation, theory packs, rendering, and zero-change checking are not external applications
but internal expressions of the same semantics. If the manuscript succeeds, the reader
should feel that the repository is building one universe of discourse for software work, not
merely a bag of clever techniques.

\paragraph{What it means to integrate the repository well.}
To integrate the repository well is not to mention every file exactly once. It is to show
how the files cooperate semantically. A good paper should therefore explain why
\repo{src/deppy/types/refinement.py} matters to \repo{src/deppy/hybrid/mixed_mode/compiler.py},
why \repo{src/deppy/hybrid/core/types.py} matters to \repo{agent/orchestrator.py}, why
\repo{src/deppy/equivalence/} and \repo{src/deppy/render/} are easier to read after one
understands gluing obstructions, and why the zero-change APIs in \repo{deppy.hybrid} are
not merely convenience wrappers but part of a deliberate adoption strategy.

\paragraph{A minimal slogan.}
If a single slogan is needed, it is this:
\begin{quote}
\emph{The plain type system gives \DepPy{} a language of exact semantic obligations.
The hybrid system gives \DepPy{} a language of layered evidence about those obligations.
Cohomology is the bookkeeping device that says whether the layers and local regions really
glue.}
\end{quote}

The remaining chapters expand that slogan until it can carry almost everything the
repository is trying to do.

The reader is therefore invited to hold one question throughout the rest of the manuscript:
\\emph{if this really is one project, what stays invariant as the objects get richer?} The
answer offered by the following chapters is that locality, transport, gluing, obstruction,
repair, and explicit evidence remain invariant. Everything else is the elaboration demanded
by software reality.
